%======================================================================
\section{What Can Data Fix?}
%======================================================================
Models B, C, D and the optional D+B are trained with exactly the protocol of model A
(\path{q2/train.py}, configurations \texttt{B}, \texttt{C}, \texttt{D}, \texttt{DB}). B draws
$g\in G$ uniformly for every training window and applies it to input and target. C, D and D+B draw
$k\in\{1,2,3\}$ uniformly and then the start frame uniformly from the allowed starts for that $k$; D+B
also draws $g$. Training videos with only two existing targets contribute equally, through the masked loss.

\qa{11} \emph{Included exactly as first written. It was committed (\texttt{Question2/predictions/Q2\_11.md},
commit \texttt{20f3c69}, 2026-10-29 04:28 PKT) after the Section~3 results and before any of B--D was trained
(B seed 0 started at 04:29 PKT).}

\begin{quote}\small
Model B (symmetry augmentation) will shrink C\_M and C\_R several-fold and bring E\_M and
E\_R down to its E(5) at 1x, but will not change the speed failure. Model C (speed augmentation) will
largely fix 2x, where the motion (about 6.6 px per step) is still within model A's 8 px reach, but only
partly fix 3x (about 10 px per step, beyond that reach), and will not change the consistency gaps.
Model D (reach 16 px plus speed augmentation) will also fix 3x, bringing E(2) at 3x close to its value at
2x at a small cost in E(5) at 1x from the coarser 8x8 grid, while its consistency gaps stay nonzero like
those of A and C.
\end{quote}

\begin{table}[H]
\centering
\caption{Models A--D and the optional D+B on the test set. $C_T$ and $E_T$ at 1$\times$; $\Ex2$ uses the
first two horizons. Training time is the mean wall-clock time per epoch over 30 epochs; inference time is
for all 1,000 test videos at batch 32 (hardware in the introduction to this question).}
\label{tab:q2-t5}
\resizebox{\linewidth}{!}{\begin{tabular}{lccccc}
\toprule
 & A & B & C & D & D+B \\
\midrule
$E^{(5)}$ 1$\times$ & 0.0224 (0.0228, 0.0219) & 0.0224 (0.0232, 0.0215) & 0.0271 (0.0276, 0.0266) & 0.0258 (0.0260, 0.0257) & 0.0266 (0.0265, 0.0267) \\
$C_M$ & 0.00156 (0.00148, 0.00164) & 0.00133 (0.00140, 0.00126) & $9.85\times10^{-4}$ ($7.42\times10^{-4}$, 0.00123) & 0.00117 (0.00104, 0.00131) & 0.00137 (0.00168, 0.00106) \\
$E_M$ & 0.0230 (0.0234, 0.0226) & 0.0224 (0.0232, 0.0216) & 0.0273 (0.0278, 0.0269) & 0.0261 (0.0262, 0.0260) & 0.0266 (0.0266, 0.0267) \\
$C_R$ & 0.00275 (0.00243, 0.00307) & 0.00175 (0.00187, 0.00164) & 0.00230 (0.00161, 0.00300) & 0.00253 (0.00239, 0.00268) & 0.00145 (0.00162, 0.00128) \\
$E_R$ & 0.0240 (0.0243, 0.0238) & 0.0224 (0.0231, 0.0217) & 0.0283 (0.0286, 0.0280) & 0.0273 (0.0274, 0.0272) & 0.0266 (0.0265, 0.0267) \\
$E^{(2)}$ 1$\times$ & 0.0189 (0.0193, 0.0185) & 0.0186 (0.0191, 0.0181) & 0.0233 (0.0236, 0.0229) & 0.0236 (0.0237, 0.0236) & 0.0243 (0.0240, 0.0245) \\
$E^{(2)}$ 2$\times$ & 0.0439 (0.0428, 0.0450) & 0.0453 (0.0448, 0.0458) & 0.0291 (0.0294, 0.0288) & 0.0260 (0.0260, 0.0260) & 0.0265 (0.0265, 0.0265) \\
$E^{(2)}$ 3$\times$ & 0.0513 (0.0505, 0.0522) & 0.0525 (0.0516, 0.0534) & 0.0336 (0.0338, 0.0335) & 0.0292 (0.0292, 0.0291) & 0.0295 (0.0295, 0.0296) \\
Trainable parameters & 480,784 & 480,784 & 480,784 & 594,496 & 594,496 \\
Training time (s/epoch) & 207 (210, 203) & 204 (204, 205) & 209 (216, 202) & 59.7 (60.0, 59.4) & 59.7 (59.8, 59.5) \\
Inference time (s, 1000 videos) & 7.83 (7.83, 7.83) & 7.92 (7.96, 7.88) & 7.96 (7.96, 7.96) & 2.21 (2.20, 2.21) & 2.19 (2.18, 2.19) \\
\bottomrule
\end{tabular}
}
\end{table}

\qa{12} \emph{Model D} is model A with a coarser fold: \texttt{pixel\_unshuffle}$(x,8)$ turns a
$1\times64\times64$ frame into $64\times8\times8$. Everything else is unchanged: two peephole-free
ConvLSTM layers with 64 hidden channels and $3\times3$ kernels (stride 1, padding 1), a $1\times1$
convolution from 64 to 64 channels, \texttt{pixel\_shuffle}$(\cdot,8)$, a sigmoid, the same 9-step
closed-loop rollout without teacher forcing, and the same training protocol (with speed augmentation, as
in its row of the model table). Its parameters are $2\times\big[(64+64)\cdot4\cdot64\cdot9+256\big]$ for the
two ConvLSTM layers plus $64\cdot64+64$ for the output: \textbf{594,496}, i.e.\ $1.24\times$ model A and
under the 721,176 limit.
\emph{Reach.} The same path argument as in Q2.3 applies, but one cell of the $8\times8$ grid is now
8 original pixels. From $u_\tau$, or from $h^{(1)}_{\tau-1}$, a path crosses layer~1's and layer~2's
$3\times3$ convolutions: 2 cells $=16$\,px. From $h^{(2)}_{\tau-1}$ it crosses only layer~2's: 8\,px.
The $1\times1$ convolution and the (un)shuffles move nothing. The reach is therefore \textbf{16\,px per
step} along a row or column ($16\sqrt2\approx22.6$\,px diagonally), strictly larger than model A's 8\,px.
It covers the 3$\times$ median of 10\,px/frame and its 95th percentile of 13\,px.
\emph{Why this design.} A bigger patch doubles the reach of every convolution without adding layers, and
the cost per step goes down (an $8\times8$ grid instead of $16\times16$). A dilation of 2 would give the
same 16\,px at A's parameter count, but each hop could then reach only even cell offsets, leaving gaps in
the neighbourhood. $5\times5$ kernels at 64 channels would exceed the parameter limit. The price of the
coarser fold is spatial detail: each state cell must now represent and decode an $8\times8$ patch from 64
channels, which may cost sharpness at 1$\times$. The design was fixed a priori from this reach calculation.
No candidate designs were compared, so neither the validation nor the test split influenced it.
\emph{Initialisation.} With PyTorch's default initialisation (the one model A uses), seed~0 of model D
never left a constant output: from the first epoch it predicted the same faint grey value (about 0.05) at every
pixel, whatever the input, and its validation loss stayed at 0.03947 for all 30 epochs. Seed~1
trained normally (validation loss 0.0267 after 30 epochs). I therefore changed model D's initialisation,
for both seeds and for D+B: Xavier-uniform weights in every convolution, zero biases, and forget-gate
biases of 1 in both ConvLSTM cells, so that cell state and gradients pass through the recurrence from
the first step (\texttt{xavier\_forget1\_} in \texttt{q2/models.py}). All four D and D+B runs were then
retrained from scratch with seeds 0 and 1; every D and D+B number in this report comes from them. The decision used
only the training curve of the failed run, not the test split. The training logs of the first runs are in \path{Question2/logs/old_init/}. With the $8\times8$ grid, an epoch also takes
about 60\,s instead of about 205\,s (Table~\ref{tab:q2-t5}): the convolutions run on 4 times fewer
positions, with only slightly more input channels per position in layer~1.

\qa{13} The three failures from Testing the Baseline get three different labels.

\emph{The consistency gaps under $M$ and $R$: appearance.} Showing the model transformed digits fixes them,
and showing it more speeds does not. Model B (symmetry augmentation, nothing else changed) brings both
transformed errors down to its ordinary error: $E_M=E_R=\Ex5=0.0224$ at 1$\times$, against
$E_M=0.0230$ and $E_R=0.0240$ for A. Its $C_R$ drops by a third ($0.00275\to0.00175$) and $C_M$ by a sixth
($0.00156\to0.00133$). Speed augmentation does not do this: for C and D, $E_R$ stays 0.0012--0.0015 above $\Ex5$,
as for A (0.0016), and $C_R$ stays at 0.0023--0.0025. Adding symmetry augmentation to D (D+B) closes the gap
again: $E_R=E_M=\Ex5=0.0266$ and $C_R=0.00145$. By Q2.1--Q2.2, every mirrored or rotated motion already
occurs in training, so what B adds is mainly how transformed digits look: the label is appearance.

\emph{The failure at 2$\times$: unfamiliar motion.} Model C has model A's architecture and reach and
differs only in its training speeds. It lowers $\Ex2$ at 2$\times$ from 0.0439 to 0.0291, much closer to
its own 0.0233 at 1$\times$ than to A's failure. The 2$\times$ motion is within A's 8\,px reach, so
the only thing that changed is that the model has now seen it. In the Figure~5 video, model C moves the
digits 4.6\,px at $h=1$ instead of model A's 1.8--2.5\,px (true 7.3--7.6\,px).

\emph{The failure at 3$\times$: mostly unfamiliar motion, partly architectural reach.} Model C also lowers
$\Ex2$ at 3$\times$, from 0.0513 to 0.0336. So most of the 3$\times$ failure was unfamiliar motion too:
model C does not ``still do badly'' at 3$\times$. Model D changes only the reach relative to C and lowers
it further to 0.0292, and it also lowers 2$\times$ to 0.0260. Relative to its own 1$\times$ error, model D
loses 24\,\% at 3$\times$ ($0.0236\to0.0292$) where C loses 44\,\% ($0.0233\to0.0336$). The part of the
3$\times$ failure that remains in C is therefore architectural reach. More data cannot carry information
further: the reach is a property of the computation graph. A predicted pixel at step $\tau+1$ is a
function of only the input pixels within 8\,px (for A and C), whatever the weights are, and training data
can only change the weights. A digit that moves 10\,px per frame arrives where nothing in the previous frame
could tell the model it was coming. C can only partly work around this: a digit is about 20\,px wide, so
most of its new position still lies within 8\,px of its old pixels, and the recurrent state already
holds its velocity from earlier frames. C can therefore smear the digit in the direction of motion, but it
cannot place the leading edge exactly. The positions in the Figure~5 video show both effects. At 2$\times$,
$h=2$, model D places the two digits 0.5 and 1.3\,px from the truth, while C is 3.5 and 7.7\,px off and
its blobs fade (45 and 61 pixels above 0.2, against 83--86 in the truth). At 3$\times$, $h=1$, both C and D move the lower digit 8.0--8.4\,px (true 11.2\,px), and both miss the
upper digit's bounce off the top wall; by $h=2$ C merges the two digits into one blob by $h=2$, whereas D keeps them apart.
One caveat: D also coarsens the grid, so the reach is not the only thing that differs between C and D.
D is also better than C at 1$\times$ (0.0258 vs 0.0271), where reach does not limit either model.

\qa{14} Yes. Speed augmentation costs accuracy on the original 1$\times$ clips, and symmetry
augmentation does not. $\Ex5$ at 1$\times$ is 0.0224 for A and 0.0224 for B, but 0.0271 for C
(+21\,\%), 0.0258 for D (+15\,\%) and 0.0266 for D+B (+19\,\%), consistently in both seeds; $\Ex2$ at
1$\times$ rises from 0.0189 to 0.0233--0.0243. The cost comes from the speed augmentation, not from the
coarser grid of D: D is better than C at 1$\times$. Only a third of the training windows are now at 1$\times$,
and predictions at 1$\times$ are visibly blurrier (compare the top rows of Figure~\ref{fig:q2-f5}, model A, and
Figure~\ref{fig:q2-f6}b, the final model). A model that must cover
three speeds also has more ways to hedge between them. Averaged over the three speeds, $\Ex2$ falls from
0.0380 (A) to 0.0263 (D), a 31\,\% improvement, and at 2$\times$ and 3$\times$ model A is no better than
predicting a blank frame. If the forecaster only ever sees clips at the original speed, the change is not
worth it: A or B is the better model. As a forecaster of motion in general, which is the point of this
question, it is worth it: D trades a 15\,\% loss at 1$\times$ for a 41--43\,\% gain at 2$\times$ and
3$\times$.

\emph{How the Q2.11 prediction fared.} Mostly right, with two misses. B did bring $E_M$ and $E_R$ down to its
$\Ex5$ at 1$\times$ and did not change the speed failure, but it shrank $C_M$ and $C_R$ by 15\,\% and
36\,\%, not ``several-fold''. C largely fixed 2$\times$ and partly fixed 3$\times$, as predicted; its
consistency gaps stayed of the same size as A's ($C_M$ even fell somewhat). D brought $\Ex2$ at 3$\times$
(0.0292) close to 2$\times$ (0.0260), with gaps still nonzero, as predicted. But its cost at 1$\times$ came
from the speed augmentation it shares with C, not from the coarser grid.

%======================================================================
\section{The Final Model}
%======================================================================
\emph{The symmetrised forecaster is exactly consistent.} Let $T\in G$, and recall that every $g\in G$
acts on frames as a pixel permutation, so it is linear and commutes with averaging. Then
\[
\bar f(TX)=\frac18\sum_{g\in G} g^{-1} f(gTX)
=\frac18\sum_{g'\in G}T\,g'^{-1} f(g'X)=T\Big(\frac18\sum_{g'\in G}g'^{-1}f(g'X)\Big)=T\bar f(X),
\]
substituting $g'=gT$. As $g$ runs over $G$, so does $g'$, because $G$ is a group; and
$g^{-1}=(g'T^{-1})^{-1}=Tg'^{-1}$. Hence $\bar f(T(X))=T(\bar f(X))$ for every input, and
$C_T=\mathrm{MSE}\big(\bar f(T(X)),T(\bar f(X))\big)=0$ for $T=M$ and $T=R$ (and every other element of
$G$), whatever $f$ is. In floating point, the two sides add the same eight terms in a different order,
so the measured gaps are rounding errors only (Table~\ref{tab:q2-t6}). The implementation is
\texttt{Symmetrised} in \texttt{q2/models.py}; inference timing includes the eight transformations, the
eight forward passes, the inverse transformations and the averaging.

\emph{Choice of final model: symmetrised D.} The default is model D, and Table~\ref{tab:q2-t5} does not
make a case for D+B. Symmetry augmentation on top of D buys learned consistency ($C_R$ $0.00253\to0.00145$,
$E_R$ $0.0273\to0.0266$ at 1$\times$), but it costs ordinary accuracy at 1$\times$ ($\Ex5$
$0.0258\to0.0266$, in both seeds) and gains nothing at 2$\times$ or 3$\times$ (0.0260 vs 0.0265,
0.0292 vs 0.0295). Once wrapped, any model is exactly consistent, so what D+B adds is redundant.
Table~\ref{tab:q2-t6} confirms this: $\bar f$(D) and $\bar f$(D+B) differ by at most 0.0002 in every
error column, with $\bar f$(D) slightly ahead at 1$\times$ (0.0258 vs 0.0260), at the same parameter count and cost.

\begin{table}[H]
\centering
\caption{Symmetrised forecasters $\bar f$ built from the seed-0 and seed-1 models B, D and D+B. Training time is
that of the underlying learned model (symmetrisation needs no training); inference time includes all
eight transformed forward passes.}
\label{tab:q2-t6}
\resizebox{\linewidth}{!}{\begin{tabular}{lccc}
\toprule
 & $\bar f$(B) & $\bar f$(D) & $\bar f$(D+B) \\
\midrule
$E^{(5)}$ 1$\times$ & 0.0217 (0.0224, 0.0209) & 0.0258 (0.0260, 0.0256) & 0.0260 (0.0258, 0.0262) \\
$C_M$ & $9.37\times10^{-17}$ ($9.44\times10^{-17}$, $9.31\times10^{-17}$) & $6.41\times10^{-17}$ ($5.82\times10^{-17}$, $7.00\times10^{-17}$) & $6.62\times10^{-17}$ ($7.21\times10^{-17}$, $6.02\times10^{-17}$) \\
$E_M$ & 0.0217 (0.0224, 0.0209) & 0.0258 (0.0260, 0.0256) & 0.0260 (0.0258, 0.0262) \\
$C_R$ & $7.55\times10^{-17}$ ($7.61\times10^{-17}$, $7.49\times10^{-17}$) & $5.13\times10^{-17}$ ($4.66\times10^{-17}$, $5.61\times10^{-17}$) & $5.30\times10^{-17}$ ($5.82\times10^{-17}$, $4.79\times10^{-17}$) \\
$E_R$ & 0.0217 (0.0224, 0.0209) & 0.0258 (0.0260, 0.0256) & 0.0260 (0.0258, 0.0262) \\
$E^{(2)}$ 1$\times$ & 0.0181 (0.0185, 0.0177) & 0.0237 (0.0238, 0.0236) & 0.0239 (0.0237, 0.0241) \\
$E^{(2)}$ 2$\times$ & 0.0444 (0.0439, 0.0449) & 0.0260 (0.0260, 0.0260) & 0.0261 (0.0261, 0.0261) \\
$E^{(2)}$ 3$\times$ & 0.0512 (0.0504, 0.0520) & 0.0289 (0.0291, 0.0288) & 0.0290 (0.0290, 0.0291) \\
Trainable parameters & 480,784 & 594,496 & 594,496 \\
Training time (s/epoch) & 204 (204, 205) & 59.7 (60.0, 59.4) & 59.7 (59.8, 59.5) \\
Inference time (s, 1000 videos) & 64.3 (64.4, 64.2) & 17.7 (17.8, 17.6) & 17.8 (17.7, 17.8) \\
\bottomrule
\end{tabular}
}
\end{table}

\begin{figure}[H]
\centering
\includegraphics[width=0.92\linewidth]{fig6a_final_M_9269.png}\\[2pt]
\includegraphics[width=0.92\linewidth]{fig6a_final_M_9256.png}
\caption[Final model, Figure 3 and Figure 5 repeated]{(a) Figure~\ref{fig:q2-f3} repeated with the final
model (symmetrised D, seed 0) on the same videos: $T=M$, videos 9269 (top) and 9256 (bottom); continued
below. Rows as in Figure~\ref{fig:q2-f3}. GIFs: \path{Question2/gifs/fig6a_final_M_video9269.gif},
\path{Question2/gifs/fig6a_final_M_video9256.gif}, \path{Question2/gifs/fig6a_final_R_video9983.gif},
\path{Question2/gifs/fig6a_final_R_video9733.gif}.}
\label{fig:q2-f6}
\end{figure}
\begin{figure}[H]
\ContinuedFloat
\centering
\includegraphics[width=0.92\linewidth]{fig6a_final_R_9983.png}\\[2pt]
\includegraphics[width=0.92\linewidth]{fig6a_final_R_9733.png}
\caption{(continued) (a) $T=R$, videos 9983 (top) and 9733 (bottom). Rows 3 and 4 now coincide, because
$\bar f(R(X))=R(\bar f(X))$.}
\end{figure}
\begin{figure}[H]
\ContinuedFloat
\centering
\includegraphics[width=\linewidth]{fig6b_speed_final.png}
\caption{(continued) (b) Figure~\ref{fig:q2-f5} repeated with the final model (symmetrised D, seed 0) on test
video 9021 at 1$\times$, 2$\times$ and 3$\times$ ($t=1$).}
\end{figure}

Measured in the same way as for model A in Q2.9 (\texttt{python -m q2.positions sym-D}), the final model
follows the digits of Figure~\ref{fig:q2-f6}b. At 1$\times$ both digits are within 1.3\,px of the truth at
$h=1$ (model A: within 0.5\,px). At $h=2$ it also misses the 1's bounce off the left wall, but by less
(4.1\,px, model A 6.2\,px). At 2$\times$ it is within 1.5\,px at $h=1$ and within 1.8\,px at $h=2$,
where model A was 4.9--8.8\,px off. At 3$\times$, $h=1$, it moves the lower digit 7.8\,px of the true
11.2\,px (4.7\,px off) and the upper digit, which bounces off the top wall, ends 3.2\,px off. Its blobs
are larger than A's at 1$\times$, but smaller at 2$\times$ and 3$\times$.

\qa{15} \emph{Consistency.} Built-in symmetry is exact and learned symmetry is not. Symmetrised B has
$C_M=9.4\times10^{-17}$ and $C_R=7.6\times10^{-17}$, rounding error in single precision, against
$C_M=0.00133$ and $C_R=0.00175$ for model B itself, which is only somewhat lower than model A. The
argument above shows why: the group average makes $\bar f(T(X))=T(\bar f(X))$ an identity that holds for
every input and every set of weights. Augmentation only rewards the network for being \emph{right} on
transformed clips, on average over the training data. Nothing ties its prediction for $T(X)$ to its
prediction for $X$, and with finite capacity and training it reaches neither on unseen videos.
\emph{Accuracy.} On transformed clips, $E_M=E_R=\Ex5$ for every wrapped model, by construction. On the
original clips, wrapping did not cost accuracy anywhere: $\Ex5$ at 1$\times$ goes from 0.0224 to 0.0217
for B ($-3\,\%$, both seeds), stays at 0.0258 for D, and goes from 0.0266 to 0.0260 for D+B. For D, the
transformed errors fall from 0.0261 ($M$) and 0.0273 ($R$) to 0.0258. $\Ex2$ at 2$\times$ and 3$\times$
changes by at most 0.0009.
\emph{Inference cost.} Wrapping costs eight forward passes: 7.92\,s to 64.3\,s for B and 2.21\,s to 17.7\,s
for D on the 1,000 test videos (about $8\times$), with no training.
\emph{Blur.} The average of eight predictions is blurred wherever they disagree. Because MSE is convex,
the wrapped forecaster's error is never larger than the mean of the eight transformed errors $E_g$ of
$f$. It can, however, be larger than $f$'s own error on the original clips, if $f$ is much better on
untransformed clips than on transformed ones, i.e.\ when $E_T\gg\Ex5$. That is the case for a model trained without
symmetry augmentation whose motion rule depends on appearance (model A, $E_R-\Ex5=0.0016$), and for
videos where the eight views give different positions. Here it did not hurt: the gaps $E_T-\Ex5$ of B
and D are small, and $\Ex5$ at 1$\times$ improved or stayed the same, so the blurring averages out
errors more than it smears correct predictions. The largest change for the worse in any column is $+0.0001$ ($\Ex2$ at
1$\times$ for D).
\qa{16} Only partly. In Q2.7 the failure on video 9983 ($T=R$) had two parts: the two accounts of
the motion disagreed (rows 3 and 4 of Figure~\ref{fig:q2-f3}c), and neither one separated the overlapping
3 and 0. The first part is fixed. In Figure~\ref{fig:q2-f6}a (video 9983), rows 3 and 4 are identical, and
the video's $C_R$ is 0 up to rounding (model A, seed 0: 0.00502). The change that fixed it is the symmetrisation; symmetry
augmentation alone (model B) only reduced it to 0.00216. The second part is not fixed. The final model still
predicts one tall blob that drifts slightly downwards and fades, while the true digits separate and move
to the lower left. On this video it is no more accurate than model A: its error is 0.0338 on both the
original and the rotated clip, against 0.0292 and 0.0335 for model A. The speed augmentation made this
1$\times$ video worse: model C gives 0.0361 and model D 0.0330 on the original clip. This matches Q2.14.
None of the changes addresses what makes this video hard: two overlapping, curly digits whose appearance
does not tell the model where each one is.
\qa{17} Partly. The final model learned a motion rule that transfers across speeds and, by construction,
across all eight symmetries. It is not a persistence or blank predictor: at 3$\times$ its $\Ex2$ is 0.0289,
against 0.0703 for Copy-last and 0.0426 for Blank, and at 2$\times$ its digits move almost the true
distance (Figure~\ref{fig:q2-f6}b). To convince a skeptic, I would measure centroid tracking errors over
the whole test set, per digit, horizon and speed, including wall bounces, against Copy-last and
constant-velocity extrapolation. I would also test an unseen speed such as 4$\times$, where a learned rule
should degrade gracefully instead of stalling as model A did at 3$\times$. One test it has not passed is
keeping digits intact: overlapping digits merge into one blob (video 9983), bounces are missed at
3$\times$, and by $h=5$ every prediction is a blurred shape rather than a digit.
